\subsection{Selected MATH Program Inventory}
\label{app:panel-mechanisms}
Table~\ref{tab:panel-mechanisms} describes the nine frozen members
selected for the repeat study. H1--H8 are neutral labels for generated
members, all constructed under seed $0$. The listed operations describe
what each program implements, not whether every branch fires or improves
correctness. The mean-call eligibility ceiling is not a per-execution
maximum, and one logical call may request several samples.

\begin{table}[htbp]
\centering\small
\begin{tabular}{@{}p{0.29\linewidth}p{0.66\linewidth}@{}}
\toprule
Frozen member & Source-defined execution path \\
\midrule
\bare{} & One temperature-$0$ solution with the final-answer format. \\
H1 (DeepSeek) & Three temperature-$0.4$ solutions; normalized-answer vote; a greedy fallback when no answer or a tie remains. \\
H2 (ERNIE) & Draft a solution, then ask the solver to verify/correct the extracted answer; return the checked answer. \\
H3 (GLM) & Up to five differently prompted temperature-$0$ solutions; stop when three agree; adjudicate unresolved disagreement, with format fallbacks. \\
H4 (Kimi) & Request five temperature-$0.7$ samples in one call; extract answers and vote. \\
H5 (Kimi) & Initial solution, verification, and an alternative solution; vote, with an additional resolution call if no majority emerges. \\
H6 (MiniMax) & Draft, critique, and an additional tie-breaking call on disagreement; parse failures use an earlier available answer. \\
H7 (MiniMax) & Request five temperature-$0.7$ samples in one call; normalize and vote, resolving ties by first occurrence. \\
H8 (Qwen) & Draft plus an independent solve-and-audit call; compare answers and request adjudication on conflict, with a format-repair path. \\
\bottomrule
\end{tabular}
\caption{\textbf{What was selected for the repeat study.} All calls use
the same frozen solver; verification and adjudication are solver prompts,
not access to benchmark labels. Six seats are builder-stratified; H5 (Kimi) and H7 (MiniMax)
are the two at-large seats.}
\label{tab:panel-mechanisms}
\end{table}
